% page 489 of the facsimile (image p-488 of the scan)
% block 177.8 x 242.0 mm ; left margin 24.2 mm
\pgc{2.540mm}{0.000mm}{
\l{}
\l{\cel{58}481}
\l{}
\l{\cel{3}raguto. (\sou{koquarto}).Karno-peco preparita per sauco. - DEFI.}
\l{}
\l{\cel{3}raitro. (olim). Kavalo-soldato de Germania, qua milito-}
\l{\cel{6}servis en Francia. - DFIRS.}
\l{}
\l{\cel{3}rajuntar. (\sou{trans.}) Irar por juntar su itere ad (ulu) de qua onu}
\l{\cel{6}esas separita. - FIS.}
\l{}
\l{\cel{3}rakar.\cel{2}(trans., de)\cel{2}Pasigar (liquido) de un vazo aden altra,}
\l{\cel{6}cherpante sube (per robineto o sifono). - E.}
\l{}
\l{\cel{3}rakahuto.\cel{2}Mixuro de fekulo, de glani dolca, de rizo-farino,}
\l{\cel{6}de tuberkuli de ter-mandeli, sukrizita ed aromatizita. - DEFIRS.}
\l{}
\l{\cel{3}raketo.\cel{2}Planketo manuo-tenebla, kun reto de tripo-kordi, uzata}
\l{\cel{6}por sendar la baloneto, en la baloneto-ludo - la flugbalono,}
\l{\cel{6}en la ludo di ta nomo. - DEFIRS.}
\l{}
\l{\cel{3}rakito. (\sou{patol.)}\cel{2}Halto di la kresko, manifestata (multa-kaze)}
\l{\cel{6}da deformeso di la spino. - DEFIRS.}
\l{}
\l{\cel{3}rakontar. (\sou{trans., ad)}. Dicar (ad ulu) narace (ulo fingita) por}
\l{\cel{6}amuzar. - EFI.}
\l{}
\l{\cel{3}ralo. (\sou{zool.)}\cel{3}Genero de stiltieri, kun grandeso mezvalora, gambi}
\l{\cel{6}pasable kurta ma tre-longa-fingra, tre apta pri kuro - vivanta}
\l{\cel{6}en la plana landi ed en la marshi ; fluganta tre poke ; cirkul-}
\l{\cel{6}anta nur dum nokto ; extreme karnavida. - L.\cel{1}\sou{rallus.}\cel{2}DEFI.}
\l{}
\l{\cel{3}raliar. (\sou{trans.})\cel{2}(\sou{milito)}.\cel{2}Asemblar sur batal-agro, la kombatanti}
\l{\cel{6}qui esis dispersita. - DEF.}
\l{}
\l{\cel{3}\sou{ramo}. (l) (\sou{bot.}) Mikra brancho qua naskis de brancho precipua.}
\l{\cel{6}(2) Subdividuro. (\sou{anat.}) Subdividuro di la vaskuli, di la nervi.}
\l{\cel{6}- (\sou{geol}.) Masivo qua divergas de montaro, aden altra direciono.}
\l{\cel{6}- (\sou{genea}l.)Subdividuro di brancho di familio, di arboro genealo-}
\l{\cel{6}giala. - DEFIS.}
\l{}
\l{\cel{3}\sou{ramno}. (\sou{bot.)}\cel{2}Arboreto di qua la beri (preske omna-kaze : nigra)}
\l{\cel{6}esas purgiva, e di qua la kortico donas kolorizivo flava. -}
\l{\cel{6}L.\cel{1}\sou{rhamnus}. - DEFS.}
\l{}
\l{\cel{3}ramio. (\sou{bot.)}\cel{3}Planto qua donas materio texebla, qua devenis de}
\l{\cel{6}Oriento. - L.\cel{1}\sou{boehmeria}.\cel{2}- DEFS.}
\l{}
\l{\cel{3}rampo. Sur voyo, plano qua formacas angulo obliqua kun la hori-}
\l{\cel{6}zontalo, e quan onu esas acensonta. - DEF.}
\l{}
\l{\cel{3}rano.\cel{2}(\sou{zool.})\cel{2}Batrako sen-kauda qua vivas, prefere, apud aquo}
\l{\cel{6}stagnanta, en la loki humida. - L.\cel{1}\sou{rana}.\cel{2}- IS.}
\l{}
\l{\cel{3}ranca.\cel{2}Dicesas pri korpo grasa (lardo, oleo, e c.) qua aquiris}
\l{\cel{6}odoro kelke mala, e saporo akra. - DEFIS.}
}
